\chapter{Loneliness}
\index{Loneliness}

\section*{Definition and Overview}
Loneliness is a subjective, distressing feeling of discrepancy between desired and actual social relationships. Unlike social isolation, which is an objective lack of social contact, loneliness refers to the perceived quality, depth, and intimacy of connections. It is a complex emotional state that can occur even in the presence of others and has significant implications for both mental and physical health. It is increasingly recognised not just as a social issue, but as a major public health determinant.

\section*{Epidemiology}
Loneliness is widely recognised as a major public health issue across all age groups. In many countries, research indicates that approximately one in four adults experience significant loneliness at any given time, with young adults (aged 18--25) and the elderly (aged 75 and older) experiencing the highest rates. Globally, the prevalence has increased in recent years, influenced by changing household structures, rapid urbanisation, and digital technologies that can replace face-to-face interactions. It affects individuals across all socioeconomic levels.

\section*{Aetiology and Pathophysiology}
Loneliness arises from a combination of social, psychological, and physiological factors.
\begin{itemize}
    \item \textbf{Social transitions:} moving to a new city, retirement, bereavement, relationship breakdown, or changes in physical mobility.
    \item \textbf{Psychological factors:} low self-esteem, social anxiety, depression, or attachment insecurities that impair social engagement.
    \item \textbf{Chronic stress response:} loneliness activates the hypothalamic-pituitary-adrenal (HPA) axis, elevating cortisol levels and promoting chronic low-grade inflammation.
    \item \textbf{Cardiovascular strain:} increased sympathetic nervous system activity leads to vasoconstriction, elevated blood pressure, and cardiovascular strain.
    \item \textbf{Sleep disruption:} altered sleep architecture and fragmented sleep due to perceived hypervigilance for threats in the social environment.
\end{itemize}

\section*{Clinical Features}
Loneliness is not a clinical diagnosis, but it manifests in diverse cognitive, emotional, and physical ways.
\begin{itemize}
    \item Persistent feelings of emptiness, sadness, or lack of belonging
    \item Social withdrawal or avoidance of social settings despite desiring connection
    \item Increased screen time, compulsive internet use, or superficial digital interactions
    \item Fatigue, low energy, and poor quality sleep (frequent waking)
    \item Heightened sensitivity to social rejection or perceived negative evaluation
    \item Symptoms of co-occurring anxiety or depressive disorders
\end{itemize}

\section*{Differential Diagnosis}
It is important to distinguish loneliness from other conditions and psychiatric disorders.
\begin{itemize}
    \item \textbf{Social isolation:} the objective state of having minimal contact with others, which may or may not be accompanied by subjective feelings of loneliness.
    \item \textbf{Major depressive disorder:} a clinical psychiatric condition characterised by persistent low mood, anhedonia, and worthlessness, which may cause or be caused by loneliness.
    \item \textbf{Social anxiety disorder:} an intense, persistent fear of being watched and judged by others, leading to avoidance of social interactions.
    \item \textbf{Agoraphobia:} an anxiety disorder characterised by fear of places or situations that might cause panic and helplessness, leading to isolation.
    \item \textbf{Apathy or anhedonia in neurodegenerative diseases:} such as dementia or Parkinson's disease, where social withdrawal is driven by cognitive or motor decline.
\end{itemize}

\section*{When to Seek Medical Care}
Individuals should seek medical or psychological support if feelings of loneliness are persistent, causing significant distress, or leading to symptoms of anxiety, depression, or self-harm.

\textbf{Contact your local emergency services for:}
\begin{itemize}
    \item Thoughts of self-harm or suicide, or feeling unable to keep oneself safe
    \item Severe panic attacks, chest pain, or acute breathing difficulty associated with extreme distress
    \item Inability to care for oneself or perform basic daily tasks due to severe depression
\end{itemize}

\section*{Diagnosis and Investigations}
Assessment focuses on exploring the individual's subjective experience and screening for associated health conditions.
\begin{itemize}
    \item \textbf{UCLA Loneliness Scale:} a validated 20-item questionnaire commonly used in clinical and research settings to measure subjective loneliness.
    \item \textbf{Social isolation assessment:} evaluating the size and frequency of the individual's social network.
    \item \textbf{Mental health screening:} using questionnaires such as the Patient Health Questionnaire-9 (PHQ-9) and Generalised Anxiety Disorder-7 (GAD-7) to assess for depression and anxiety.
    \item \textbf{Physical health assessment:} evaluating cardiovascular health, sleep hygiene, and inflammatory markers if chronic stress is suspected.
\end{itemize}

\section*{Management}
Interventions must be tailored to address the specific type and cause of loneliness.
\begin{itemize}
    \item \textbf{Social prescribing:} clinical referrals to community-based activities, interest groups, volunteering, or social clubs.
    \item \textbf{Cognitive behavioural therapy (CBT):} to address maladaptive social cognitions, fear of rejection, and hypervigilance to social threats.
    \item \textbf{Digital and social skills training:} programmes to help individuals build confidence and competence in establishing connections.
    \item \textbf{Community and lifestyle measures:} encouragement of regular physical exercise, mindfulness, and direct face-to-face interactions.
    \item \textbf{Addressing barriers:} providing transportation support or mobility aids for older adults to facilitate social participation.
\end{itemize}

\section*{Complications}
\begin{itemize}
    \item Increased risk of developing major depressive disorder and anxiety disorders
    \item Higher incidence of cardiovascular diseases, including coronary heart disease and stroke
    \item Accelerated cognitive decline and a higher risk of developing dementia in older adults
    \item Impaired immune function, leading to increased susceptibility to infections
    \item Increased mortality rates, comparable to well-established risk factors such as smoking and obesity
\end{itemize}

\section*{Prognosis and Outcomes}
The prognosis is highly variable and depends on the underlying causes and the individual's capacity to engage in interventions. Loneliness can be transient, resolving with changes in life circumstances (such as starting a new job or joining a group). However, chronic loneliness can become a self-reinforcing state due to social withdrawal and negative cognitive biases. Early intervention and social support significantly improve mental and physical health outcomes.

\section*{Prevention and Self-Care}
\begin{itemize}
    \item Foster and maintain existing friendships through regular contact and communication.
    \item Engage in volunteering or community groups to meet people with shared interests.
    \item Limit passive social media use and prioritise active, meaningful connections.
    \item Maintain a regular routine of physical activity, which can improve mood and open social avenues.
    \item Practice mindfulness or self-compassion to reduce self-criticism and social anxiety.
\end{itemize}

\section*{Sources and Further Reading}
The following reputable organisations provide further information on loneliness, social connection, and mental health support.
\begin{itemize}
    \item Australian Institute of Health and Welfare. \textit{Social isolation and loneliness in Australia.} https://www.aihw.gov.au/
    \item Healthdirect Australia. \textit{Loneliness and isolation: support and advice.} https://www.healthdirect.gov.au/loneliness-and-isolation
    \item Ending Loneliness Together. \textit{Research, resources, and campaigns to end loneliness.} https://endingloneliness.com.au/
    \item Beyond Blue. \textit{Managing loneliness and connecting with others.} https://www.beyondblue.org.au/
    \item NHS. \textit{Feeling lonely: advice and support for older people and adults.} https://www.nhs.uk/mental-health/advice-for-life-situations-and-events/feeling-lonely/
\end{itemize}
